\documentclass[11pt,a4paper]{article}
\usepackage[margin=27mm]{geometry}
\usepackage[T1]{fontenc}
\usepackage{lmodern}
\usepackage{amsmath,amssymb,amsthm,booktabs,array}
\usepackage{xurl}
\usepackage[colorlinks=true,linkcolor=blue,citecolor=blue,urlcolor=blue]{hyperref}
\newtheorem{theorem}{Theorem}
\newtheorem{corollary}[theorem]{Corollary}
\newcommand{\R}{\mathbb R}
\newcommand{\E}{\mathbb E}
\newcommand{\Dstar}{D^{\star}}
\title{Earliest swapping for every chain length:\\
Kamin et al. (2023), Section IV D 1, Eq. (41)}
\author{Haobo Ma\\\small ChronoAI Pte Ltd, Singapore
\and Wenlin Zhang\\\small National University of Singapore\\
\small ChronoAI Pte Ltd, Singapore}
\date{November 2, 2026}
\begin{document}
\maketitle
\begin{abstract}
The all-length earliest-swapping conjecture in Section IV D 1, Eq.~(41),
of Kamin, Shchukin, Schmidt and van Loock (2023) states that repeatedly
merging an adjacent pair with smallest maximum ready time minimizes total
dephasing. We give a self-contained potential proof for arbitrary
nonnegative real ready times and every tie choice. The optimum is the
largest ready time plus the sum of adjacent maxima minus the sum of ready
times. An exact nonnegative slack formula characterizes all optimal
histories. The pathwise inequality also answers the exponential-dephasing
question in Section IV D 5, for the same fixed-readiness swapping model:
earliest swapping simultaneously minimizes expected dephasing and maximizes
every expected exponential attenuation factor.
\end{abstract}

\section{The conjecture and its model}
For parallel entanglement distribution with deterministic instantaneous
swaps, Kamin et al.~\cite[Sec.~IV D 1, Eq.~(41), p.~023086-13]{kamin}
ask whether earliest adjacent-pair swapping minimizes dephasing for every
chain length. Their Eqs.~(38)--(39) define the greedy and optimal costs;
they verify equality for lengths up to eight. The natural-number result
has the pinned formal source~\cite{lean}. We give its potential argument
in real variables, and derive the exact slack and probabilistic
consequences. The lower bound is preserved by every greedy step.

Let $x=(x_1,\ldots,x_n)$, $n\geq1$, be nonnegative real ready times.
A merge at $i\in\{1,\ldots,n-1\}$ costs
$c_i(x)=|x_i-x_{i+1}|$ and replaces the adjacent entries by their maximum:
\[
 T_i x=(x_1,\ldots,x_{i-1},\max(x_i,x_{i+1}),x_{i+2},\ldots,x_n).
\]
A complete history $H$ has $n-1$ valid merges, and $C_H(x)$ is the sum
of their costs. Define $\Dstar(x)=\min_H C_H(x)$, with cost zero on a
singleton. This is exactly the Bellman recurrence in Eq.~(39).
A greedy history chooses a pair minimizing $\max(x_i,x_{i+1})$ at each
stage, with arbitrary choices among ties, as in Eq.~(38).

The readiness vector is fixed while histories are compared. The claim is
about this recurrence, rather than an optimization over distribution
schedules, memory cutoffs, or probabilistic swapping.

\section{A potential and an exact slack formula}
Write $M(x)=\max_j x_j$ and set
\begin{equation}\label{eq:potential}
 P(x)=M(x)+\sum_{j=1}^{n-1}\max(x_j,x_{j+1})-\sum_{j=1}^n x_j.
\end{equation}
For a singleton, $P=0$. The identity
$2\max(u,v)=u+v+|u-v|$ gives the equivalent expression
\begin{equation}\label{eq:variation}
 2P(x)=2M(x)-x_1-x_n+\sum_{j=1}^{n-1}|x_j-x_{j+1}|.
\end{equation}
In particular, $P(x)\geq0$.

\begin{theorem}\label{thm:main}
For every $n\geq1$ and $x\in\R_{\geq0}^n$,
\[
 \Dstar(x)=P(x).
\]
Every complete greedy history attains this value, independently of all
tie choices. More generally, if $x^{(0)}=x$ and
$x^{(s+1)}=T_{i_s}x^{(s)}$ describe any complete history, then
\begin{equation}\label{eq:regret}
 C_H(x)-P(x)=\sum_{s=0}^{n-2}\delta_{i_s}(x^{(s)}),
 \qquad \delta_i(x)\geq0,
\end{equation}
where, writing $a=x_i$, $b=x_{i+1}$ and $m=\max(a,b)$,
\begin{align}\label{eq:slack}
\delta_i(x)={}&
\begin{cases}
 \max(x_{i-1},m)-\max(x_{i-1},a),& i>1,\\
 0,&i=1
\end{cases}\notag\\
&+\begin{cases}
 \max(m,x_{i+2})-\max(b,x_{i+2}),&i+1<n,\\
 0,&i+1=n.
\end{cases}
\end{align}
Consequently, a history is optimal if and only if every one of its
merge slacks vanishes.
\end{theorem}

\begin{proof}
A merge preserves the global maximum. In the adjacent-maxima sum, it
removes the term $m$, changes the left neighboring term from
$\max(x_{i-1},a)$ to $\max(x_{i-1},m)$ when present, and changes the right
neighboring term from $\max(b,x_{i+2})$ to $\max(m,x_{i+2})$ when present.
The sum of entries changes by $m-a-b$. Hence
\[
 P(T_i x)-P(x)=\delta_i(x)+a+b-2m
             =\delta_i(x)-|a-b|.
\]
This proves $c_i(x)+P(T_i x)-P(x)=\delta_i(x)$, including merges
at either end of the list. Both summands in \eqref{eq:slack} are
nonnegative because $m\geq a,b$.

If $i$ is greedy and the left neighbor exists, then
$m\leq\max(x_{i-1},a)$ by comparison with the adjacent pair $(i-1,i)$.
It follows that
$\max(x_{i-1},m)=\max(x_{i-1},a)$.
The right comparison gives
$\max(m,x_{i+2})=\max(b,x_{i+2})$ when that neighbor exists.
Thus every greedy step has zero slack. These arguments use weak
inequalities, so ties cause no exception.

Summing the one-step identity over a complete history telescopes the
potential to its singleton value zero, proving \eqref{eq:regret}.
It follows that every history costs at least $P(x)$, while any greedy
history costs exactly $P(x)$. A greedy pair always exists in a finite
non-singleton list, and each merge decreases its length by one, so such a
complete history exists. This proves the optimum formula. Since a finite
sum of nonnegative numbers is zero exactly when all its summands are
zero, it also proves the last assertion.
\end{proof}

Greediness is sufficient but not necessary for optimality. For
$x=(1,2,4,3)$, merging the last pair first is not greedy, but costs
$1+\Dstar(1,2,4)=1+3=4=P(x)$. For monotone ready times,
\eqref{eq:variation} reduces to $P(x)=\max_jx_j-\min_jx_j$, recovering
the observation following Eq.~(40) of~\cite{kamin} for every length.

Greedy histories are also chronological: after a merge with maximum $m$,
all unchanged adjacent-pair maxima are at least $m$ by greediness, and
every newly formed adjacent pair contains the entry $m$. The next merge
time is therefore at least $m$. A merged block is ready at the maximum
of its elementary ready times, by induction. Thus these histories respect
readiness, including consecutive instantaneous swaps at a tied time.

Goodenough, Coopmans and Towsley~\cite[Sec.~3.1]{goodenough} express
swap-ASAP memory noise through the roughness
$\sum_j|x_j-x_{j+1}|$ and the endpoint waiting contribution
$2M(x)-x_1-x_n$. Their combined exponent is exactly $2P(x)$ by
\eqref{eq:variation}. The factor of two records the conversion to the
normalization of Eqs.~(38)--(39); the roughness alone is a different
objective. The potential proof above establishes optimality against
every history in those recurrences.

\section{Probabilistic consequences and source coverage}
\begin{corollary}\label{cor:attenuation}
Let $X$ have any joint distribution on $\R_{\geq0}^n$. Compare histories
using the same realized readiness vector, and allow a competing policy
to use additional randomness. For every $\alpha\geq0$,
\[
 C_{\mathrm{greedy}}(X)=P(X)\leq C_H(X),\qquad
 \E e^{-\alpha C_{\mathrm{greedy}}(X)}
 \geq\E e^{-\alpha C_H(X)}.
\]
Whenever the costs are integrable, greedy histories also minimize
their expectation.
\end{corollary}
\begin{proof}
Theorem~\ref{thm:main} applies to every realization and every history.
Apply the decreasing function $t\mapsto e^{-\alpha t}$ and then take
expectations. These exponential variables lie in $[0,1]$, so no moment
condition is needed. Taking expectations directly gives the cost
comparison when integrable.
\end{proof}

This answers the question about minimizing $\E D$ versus maximizing
$\E e^{-\alpha D}$ raised in Section IV D 5 of~\cite[p.~023086-16]{kamin},
within the parallel-distribution recurrence: the stronger, pointwise
inequality supplies both conclusions. No independence or identical
distribution assumption is needed.

\begin{corollary}\label{cor:mean}
If $X_1,\ldots,X_n$ are independent geometric random variables on
$\{1,2,\ldots\}$ with success probability $p\in(0,1]$ and $q=1-p$, then
\begin{equation}\label{eq:mean}
 \E\Dstar(X)=
 \sum_{j=1}^{n}\frac{(-1)^{j+1}\binom nj}{1-q^j}
 +\frac{n-2}{p}-\frac{n-1}{p(2-p)}.
\end{equation}
\end{corollary}
\begin{proof}
The tail-sum formula and the binomial theorem give
\[
 \E M(X)=\sum_{k=0}^{\infty}\bigl[1-(1-q^k)^n\bigr]
 =\sum_{j=1}^n\frac{(-1)^{j+1}\binom nj}{1-q^j}.
\]
Only finitely many $j$ occur, and each geometric series converges.
For two independent coordinates,
$\E\min(X_1,X_2)=\sum_{k\geq0}q^{2k}=1/[p(2-p)]$ and
$\E\max(X_1,X_2)=2/p-1/[p(2-p)]$.
Take expectations in \eqref{eq:potential} and substitute these
identities. This gives \eqref{eq:mean}, including $n=1$ and $p=1$.
\end{proof}

\begin{center}
\small
\begin{tabular}{@{}p{0.40\linewidth}p{0.54\linewidth}@{}}
\toprule
Source item in Kamin et al. (2023)&Coverage\\
\midrule
Sec.~IV D 1, Eq.~(41), all lengths and arbitrary ties &
Theorem~\ref{thm:main}; nonnegative real times included.\\[3pt]
Sec.~IV D 1, observation after Eq.~(40)&
Monotone-time formula from \eqref{eq:variation}.\\[3pt]
Sec.~IV D 5, p.~023086-16: mean dephasing versus exponential attenuation&
Corollary~\ref{cor:attenuation}, for the fixed-readiness recurrence.\\[3pt]
Sec.~IV D 5 and App.~E 6: optimal secret-key rate across distribution and measurement schemes&
Not implied by the recurrence: those comparisons change the readiness law,
storage convention, or rate objective.\\[3pt]
Memory-cutoff and probabilistic-swapping extensions&
Outside the deterministic, fixed-readiness model of Theorem~\ref{thm:main}.\\
\bottomrule
\end{tabular}
\end{center}

\section{Formal statement and reproducible finite checks}
The pinned Lean source~\cite{lean} contains the declaration
\nolinkurl{D5.S3.Quantum.Optimization.SwapEarliestDephasingOptimality.result}.
It proves the natural-number statement: every complete history satisfying
\texttt{IsGreedyRun} has \texttt{runCost = dstar}, where \texttt{dstar}
implements Eq.~(39) by length-many Bellman iterations. Empty and singleton
lists have cost zero. The source's private lemmas
\texttt{dstar\_potential} and \texttt{greedy\_cost\_potential} identify
both costs with the integer-valued potential corresponding to
\eqref{eq:potential}. The direct file link is
\href{https://github.com/the-omega-institute/trureturing/blob/83b70a1d0d6a58bb32521c741805d845d401c5f8/D5/S3/Quantum/Optimization/SwapEarliestDephasingOptimality.lean}{\texttt{SwapEarliestDephasingOptimality.lean}}, at commit
\nolinkurl{83b70a1d0d6a58bb32521c741805d845d401c5f8}.
The real-time extension, exact slack characterization, and probabilistic
corollaries have the written proofs given above.

The accompanying standard-library Python program \texttt{verify.py}
compares an exhaustive Bellman recursion with the smallest and largest
costs among all greedy tie continuations. It checks every vector over
$\{0,1,2,3\}$ through length eight, every vector over
$\{0,\tfrac13,\tfrac12,2\}$ through length six, and all $8!$ permutations
of eight distinct readiness ranks. Every valid first merge also satisfies
the slack identity, its nonnegativity, and the optimal-first-step
characterization. Additional exact rational comparisons check
\eqref{eq:mean} for $n=1,\ldots,6$ and $p\in\{\tfrac13,\tfrac12,1\}$
against tail sums with an explicit geometric remainder bound.
These finite checks complement the general written proof.

\begin{thebibliography}{9}
\bibitem{kamin}
L. Kamin, E. Shchukin, F. Schmidt and P. van Loock,
\emph{Exact rate analysis for quantum repeaters with imperfect memories
and entanglement swapping as soon as possible},
Physical Review Research \textbf{5}, 023086 (2023).
\href{https://doi.org/10.1103/PhysRevResearch.5.023086}{doi:10.1103/PhysRevResearch.5.023086}.
\href{https://arxiv.org/abs/2203.10318v1}{arXiv:2203.10318v1} (2022).

\bibitem{goodenough}
K. Goodenough, T. Coopmans and D. Towsley,
\emph{On noise in swap ASAP repeater chains: exact analytics,
distributions and tight approximations},
Quantum \textbf{9}, 1744 (2025).
\href{https://doi.org/10.22331/q-2025-05-15-1744}{doi:10.22331/q-2025-05-15-1744}.
\href{https://arxiv.org/abs/2404.07146v5}{arXiv:2404.07146v5}.

\bibitem{lean}
\texttt{trureturing} contributors,
\emph{SwapEarliestDephasingOptimality.lean},
\texttt{D5/S3/Quantum/Optimization}, commit
\texttt{83b70a1d0d6a58bb32521c741805d845d401c5f8}.
\href{https://github.com/the-omega-institute/trureturing/blob/83b70a1d0d6a58bb32521c741805d845d401c5f8/D5/S3/Quantum/Optimization/SwapEarliestDephasingOptimality.lean}{Pinned formal source}.
\end{thebibliography}

\medskip\noindent
\textbf{License.} Creative Commons Attribution 4.0 International
(CC BY 4.0; \texttt{cc-by-4.0}).
\end{document}
